\documentclass[11pt]{article}

\usepackage[margin=1in]{geometry}
\usepackage[utf8]{inputenc}
\usepackage{amsmath}
\usepackage{amssymb}
\usepackage{mathtools}
\usepackage{booktabs}
\usepackage{graphicx}
\usepackage{float}
\usepackage{hyperref}
\usepackage[capitalize,noabbrev]{cleveref}
\usepackage{xcolor}
\usepackage{textcomp}

\definecolor{darkblue}{rgb}{0.19, 0.19, 0.62}
\hypersetup{
  colorlinks = true,
  citecolor  = darkblue,
  filecolor  = darkblue,
  urlcolor   = darkblue,
}

\title{
  Improving Qwen for Mathematical Reasoning
}

\author{%
  Luis Xu, Wenhao Ren\\
  Department of Computer Science\\
  University of California San Diego\\
  \texttt{w4ren@ucsd.edu, lxxu@ucsd.edu}
}

\begin{document}
\maketitle

\begin{abstract}
We study mathematical problem solving with Qwen3 4B Thinking 2507 under the competition rule that the final response must come from this model and cannot use external tools or larger models. A direct prompting pipeline often produces a long solution but misses the exact answer format required by the judge. We focus on that gap. We compare raw Qwen generation, a LoRA answer finalizer, a candidate reranking stage, a simple problem type rule, and an exploratory GRPO adapter. On a disjoint public held out set of 40 problems, raw Qwen with a 2048 token budget scored 14/40. The LoRA finalizer raised this to 25/40. Candidate reranking also scored 25/40, with slightly better multiple choice accuracy and slightly worse free response accuracy. The problem type rule was best on an earlier 20 problem development set, but it fell to 24/40 on the held out set. GRPO did not beat supervised finalization. We have not yet submitted these results to the private Kaggle leaderboard.
\end{abstract}

\section{Introduction}
\label{sec:intro}

\paragraph{Problem Definition.}
The competition asks us to improve the math ability of Qwen3 4B Thinking 2507. Each row gives a math problem and expects either one multiple choice option or one or more free response answers. The problems cover calculation, algebra, statistics, calculus, number theory, and contest style reasoning. The final answer must be produced by the allowed Qwen model. Prompting, supervised fine tuning, and reinforcement learning are allowed. External APIs, other models, calculators, and code tools are not allowed at inference time. The score is unified accuracy, so every row counts the same.

\paragraph{Problem Significance.}
Small language models are useful only if they can be trusted on exact tasks. Math is a good stress test because small mistakes in arithmetic, symbolic manipulation, or answer formatting make the whole response wrong. Better math reasoning in a 4 billion parameter model would make local tutoring, quantitative analysis, and STEM support cheaper to run. This project asks whether a fixed small model can be made more reliable with careful inference design and lightweight adaptation.

\paragraph{Technical Challenge.}
The hard part is not just solving the math. The model also has to give the answer in the exact form that the judge expects. Raw Qwen often writes a plausible solution but then gives the wrong final form, omits subanswers, or runs out of budget before finishing. Free response rows can contain many blanks, so a single final number is not enough. Multiple choice rows have a different issue because the model must finish with one clean option letter. These failure modes led us to separate the reasoning step from the answer formatting step.

\paragraph{Contributions.}
The starter pipeline runs Qwen and scores the result directly. We built a small set of Qwen based variants and compared them on both an early development slice and a separate held out slice. The most important result is that answer finalization gives the largest reliable gain.

Our main contributions are:
\begin{itemize}
  \item \textbf{Answer finalization.} We trained a LoRA adapter that reads the problem and a Qwen reasoning trace, then outputs only the final boxed answer. On the held out set, this improved raw Qwen from 14/40 to 25/40.
  \item \textbf{Candidate reranking.} We generated 1024 token and 2048 token candidates, finalized both answers, and used Qwen again only when the answer keys disagreed. This tied the LoRA finalizer at 25/40, but it changed which rows were correct.
  \item \textbf{A check on selection bias.} A rule that used one method for multiple choice rows and another for free response rows scored best on the first 20 rows. It did not stay best on the held out 40 rows. This was a useful warning that the first slice was too small for choosing a final pipeline.
  \item \textbf{GRPO ablation.} We trained a GRPO adapter with the local judge as an outcome reward. Direct GRPO scored 17/40. GRPO used as a finalizer matched the LoRA finalizer at 25/40 but did not improve it.
\end{itemize}

\section{Related Work}
\label{sec:related}

\paragraph{Reasoning prompts.}
Chain of thought prompting showed that intermediate reasoning can improve model performance on math and symbolic tasks \cite{wei2022cot}. Self consistency samples several reasoning paths and chooses the most common final answer \cite{wang2022selfconsistency}. Our reranking experiment follows the same basic idea that extra test time compute can help, but our results are mixed. It helps some multiple choice rows and hurts some free response rows.

\paragraph{Lightweight adaptation.}
LoRA trains low rank update matrices while keeping the base model fixed \cite{hu2021lora}. This is a natural fit here because many errors are answer formatting errors, not complete reasoning failures. We use LoRA as a final answer adapter instead of training a new solver from scratch.

\paragraph{Reinforcement learning for math.}
GRPO was used in the DeepSeekMath work as a memory efficient policy optimization method for math models \cite{shao2024deepseekmath}. We tried it because the competition judge gives a direct outcome reward. In our setup the reward was sparse, and the adapter did not improve over supervised finalization.

\section{Methods}
\label{sec:methods}

\paragraph{Problem Setting.}
Let $x_i$ be a problem and $a_i$ be the gold answer. The model writes a response $y_i$. The local judge extracts the final answer and returns $J(y_i, a_i) \in \{0,1\}$. We want to maximize
\[
    \frac{1}{N}\sum_{i=1}^{N} J(y_i, a_i).
\]
For multiple choice rows, $a_i$ is one option letter. For free response rows, $a_i$ can contain several ordered answers.

\paragraph{Baseline: Raw Qwen Reasoning.}
The baseline prompts Qwen3 4B Thinking 2507 to reason and then place the final answer in a box. We use a 2048 token completion budget and score the raw output directly. This measures what the base model can do without a separate formatting or checking step.

\paragraph{Method 1: LoRA Answer Finalizer.}
Method 1 keeps the raw Qwen reasoning trace but adds a final formatting pass. Qwen first writes a 2048 token solution. A LoRA adapter then reads the original problem and the solution trace and writes the final boxed answer. The adapter is trained on public labeled rows using next token prediction. In other words, it learns to map a problem and a rough solution into the answer string the judge expects.

\paragraph{Method 2: Candidate Reranking.}
Method 2 adds a second solution and only spends extra compute when the answers conflict. We generate one 1024 token candidate and one 2048 token candidate. The LoRA finalizer extracts an answer from each. If the two final answers disagree, Qwen compares the candidates or solves the problem again. The reranker output is then sent through the LoRA finalizer. This is meant to catch cases where one sampled path reaches the right answer and another does not.

\paragraph{Method 3: Problem Type Selection.}
Method 3 was a simple rule based on the first 20 public rows. On that slice, Method 1 was better for multiple choice rows and Method 2 was better for free response rows. We therefore used Method 1 on multiple choice rows and Method 2 on free response rows. This rule did not hold up on the held out set, so we treat it as an attempted method rather than the final choice.

\paragraph{Exploratory GRPO.}
We also trained a GRPO adapter. The policy started from the answer finalizer LoRA. The reward was the local judge score, where a correct final answer receives 1 and an incorrect answer receives 0. The run excluded the 40 held out ids from training. It used 256 examples, 60 maximum optimization steps, a 512 token prompt budget, and a 128 token completion budget. A larger profile did not fit on the available 11GB GPUs.

\begin{figure}[H]
\centering
\fbox{
\begin{minipage}{0.92\linewidth}
\scriptsize
\[
\begin{array}{rcl}
\text{Input}
& \rightarrow &
\text{Qwen 2048}
\rightarrow
\text{score raw}
\quad \text{Baseline}
\\[6pt]
\text{Input}
& \rightarrow &
\text{Qwen 2048}
\rightarrow
\text{LoRA finalizer}
\quad \text{Method 1}
\\[6pt]
\text{Input}
& \rightarrow &
\left.
\begin{array}{l}
\text{Qwen 1024} \rightarrow \text{LoRA finalizer} \\
\text{Qwen 2048} \rightarrow \text{LoRA finalizer}
\end{array}
\right\}
\rightarrow
\text{rerank}
\rightarrow
\text{LoRA finalizer}
\quad \text{Method 2}
\\[10pt]
\text{Input}
& \rightarrow &
\text{Multiple choice uses Method 1, free response uses Method 2}
\quad \text{Method 3}
\end{array}
\]
\end{minipage}
}
\caption{Pipeline overview. Method 1 adds an answer finalizer to one Qwen reasoning pass. Method 2 compares two candidate lengths when their extracted answers disagree. Method 3 applies a simple rule based on problem type.}
\label{fig:pipeline}
\end{figure}

\section{Experiments}
\label{sec:experiments}

We compare raw Qwen generation, LoRA finalization, candidate reranking, problem type selection, and GRPO. We use the local judge for all numbers.

\subsection{Baselines}
\label{sec:baselines}

We evaluate five systems:
\begin{itemize}
  \item \textbf{Method 0: Raw Qwen 2048.} Qwen writes one response with a 2048 token budget. The raw response is scored directly.
  \item \textbf{Method 1: Qwen 2048 plus LoRA finalizer.} Qwen writes a solution trace, and the LoRA finalizer writes the boxed answer.
  \item \textbf{Method 2: 1024 and 2048 reranking plus LoRA.} The system generates two candidates, finalizes both, reranks conflicts, and finalizes the reranker output.
  \item \textbf{Method 3: Problem type selection.} The system uses Method 1 for multiple choice rows and Method 2 for free response rows.
  \item \textbf{GRPO.} A GRPO adapter is tested as a direct generator and as a finalizer.
\end{itemize}

\subsection{Evaluation}
\label{sec:evaluation}

The public set contains 1126 rows. Of these, 375 are multiple choice and 751 are free response. The free response rows contain 2273 total answer slots. We used the first 20 public rows for early development. We then built a separate held out slice of 40 rows that excludes ids 0 through 19. This slice has 20 multiple choice rows and 20 free response rows. The free response rows include 9 one slot rows, 5 two slot rows, 4 rows with 3 to 4 slots, and 2 rows with at least 5 slots.

We report multiple choice accuracy, free response accuracy, overall accuracy, balanced accuracy, and a public weighted accuracy estimate. Balanced accuracy averages the two row types equally. Public weighted accuracy uses the public row mix:
\[
    0.333 \cdot \text{multiple choice accuracy}
    +
    0.667 \cdot \text{free response accuracy}.
\]

\begin{table}[H]
\caption{Dataset and validation slice statistics. The held out slice is balanced for analysis and is not meant to match the full public distribution.}
\centering
\resizebox{\linewidth}{!}{
\begin{tabular}{lrrrrl}
\toprule
Split & Total rows & MCQ & Free response & Free slots & Notes \\
\midrule
Full public set & 1126 & 375 & 751 & 2273 & Natural public distribution \\
First 20 development slice & 20 & 9 & 11 & 22 & Used for method selection \\
Held out slice & 40 & 20 & 20 & 55 & Excludes ids 0 through 19 \\
\bottomrule
\end{tabular}
}
\end{table}

\begin{table}[H]
\caption{Composition of the 40 row held out slice. The free response subset includes both short and multi answer rows.}
\centering
\begin{tabular}{lr}
\toprule
Group & Count \\
\midrule
MCQ & 20 \\
Free response, 1 answer slot & 9 \\
Free response, 2 answer slots & 5 \\
Free response, 3 to 4 answer slots & 4 \\
Free response, 5 or more answer slots & 2 \\
\midrule
Total & 40 \\
\bottomrule
\end{tabular}
\end{table}

\subsection{Implementation Details}
\label{sec:implementationdetails}

All inference uses Qwen3 4B Thinking 2507 with Hugging Face Transformers. The answer finalizer LoRA was trained for one epoch on 1112 public answer formatting examples. Fourteen examples were skipped because they were too long. We used LoRA rank 16, alpha 32, dropout 0.05, learning rate 0.0002, batch size 1, and gradient accumulation 8. The finalizer took about 6 minutes to train on one NVIDIA A30 GPU with 24GB VRAM.

The main reasoning pass uses a 2048 token budget. Method 2 also uses a 1024 token candidate. The finalizer uses a short budget because it only needs to write the final boxed answer. We chose these values after early tests showed that longer completions often spent more time reasoning without fixing answer formatting.

\subsection{Results}
\label{sec:results}

The first 20 row development results are shown below. The biggest gain came from the LoRA finalizer. Raw Qwen scored 7/20. The same reasoning traces with finalization scored 12/20. Reranking raised free response accuracy from 6/11 to 7/11 but lowered multiple choice accuracy from 6/9 to 5/9. This motivated Method 3, which scored 13/20 on the development slice.

\begin{table}[H]
\caption{Development results on the first 20 public rows. These numbers motivated Method 3, but the held out results show that the rule did not generalize.}
\centering
\begin{tabular}{lccc}
\toprule
Method & Free response & MCQ & Overall \\
\midrule
Raw Qwen 2048, strict extraction & 6/11 & 1/9 & 7/20 \\
Qwen 2048 plus LoRA finalizer & 6/11 & 6/9 & 12/20 \\
1024 and 2048 conflict rerank plus LoRA & 7/11 & 5/9 & 12/20 \\
Problem type selection & 7/11 & 6/9 & 13/20 \\
\bottomrule
\end{tabular}
\end{table}

The main held out result is shown below. Method 1 and Method 2 tie at 25/40. Method 1 is simpler and has the better public weighted score because it is stronger on free response rows. Method 3 falls to 24/40, which suggests that the rule was tuned too closely to the first 20 rows.

\begin{table}[H]
\caption{Main results on the 40 row held out slice. Balanced accuracy averages the two row types equally. Public weighted accuracy uses the public row mix.}
\centering
\resizebox{\linewidth}{!}{
\begin{tabular}{lccccc}
\toprule
Method & MCQ & Free response & Overall & Balanced accuracy & Public weighted accuracy \\
\midrule
Raw Qwen 2048 & 4/20 & 10/20 & 14/40 & 0.3500 & 0.4001 \\
Qwen 2048 plus LoRA finalizer & 11/20 & 14/20 & 25/40 & 0.6250 & 0.6501 \\
1024 and 2048 rerank plus LoRA & 12/20 & 13/20 & 25/40 & 0.6250 & 0.6334 \\
Problem type selection & 11/20 & 13/20 & 24/40 & 0.6000 & 0.6167 \\
\bottomrule
\end{tabular}
}
\end{table}

\begin{table}[H]
\caption{The first 20 rows were not enough for final method selection.}
\centering
\resizebox{\linewidth}{!}{
\begin{tabular}{lccc}
\toprule
Method & Development slice & Held out slice & Interpretation \\
\midrule
Raw Qwen 2048 & 7/20 & 14/40 & Weak baseline, especially on MCQ \\
Qwen 2048 plus LoRA finalizer & 12/20 & 25/40 & Most reliable gain \\
1024 and 2048 rerank plus LoRA & 12/20 & 25/40 & Same total, different errors \\
Problem type selection & 13/20 & 24/40 & Did not generalize \\
\bottomrule
\end{tabular}
}
\end{table}

\subsection{Ablations and Negative Results}
\label{sec:ablations}

We also tested several ideas that did not become the main pipeline. A 3072 token prompt ablation scored 6/10. A three sample 1024 token voting run also scored 6/10. These runs did not show that more tokens or simple voting were enough. GRPO was more useful than raw Qwen in direct mode, but it still trailed the LoRA finalizer by a wide margin.

\begin{table}[H]
\caption{Ablations and negative results. GRPO direct improved over raw Qwen on some rows, but it did not beat supervised finalization.}
\centering
\resizebox{\linewidth}{!}{
\begin{tabular}{llllccccc}
\toprule
Experiment & Split & Eval mode & Setup & MCQ & Free response & Overall & Balanced & Takeaway \\
\midrule
3072 token prompt & 10 row dev & Direct & Longer generation & NA & NA & 6/10 & NA & More tokens were not enough \\
Three sample voting & 10 row dev & Vote & Three samples & NA & NA & 6/10 & NA & Too little useful diversity \\
GRPO direct & Held out 40 & Direct & Outcome reward & 8/20 & 9/20 & 17/40 & 0.4250 & Below finalization \\
GRPO finalizer & Held out 40 & Finalizer & Outcome reward & 11/20 & 14/20 & 25/40 & 0.6250 & Matches Method 1 only \\
\bottomrule
\end{tabular}
}
\end{table}

\subsection{Qualitative Comparison}
\label{sec:qualitative}

The qualitative examples below show representative rows from the held out set. The examples make the error pattern clearer than the aggregate score alone. Finalization fixes many answer format issues, but it can still collapse a multi answer problem into one number. Reranking can fix a wrong multiple choice answer, but it can also replace a correct free response answer.

\begin{table}[H]
\caption{Qualitative examples from the held out set. The table gives final answers or short behavior summaries.}
\centering
\resizebox{\linewidth}{!}{
\begin{tabular}{p{0.08\linewidth}p{0.20\linewidth}p{0.19\linewidth}p{0.16\linewidth}p{0.16\linewidth}p{0.16\linewidth}}
\toprule
ID & Problem type & Raw Qwen 2048 & Method 1 & Method 2 & Takeaway \\
\midrule
20 & Free response, 15 slots & Reached the budget and gave only the final standard deviation & Incorrect: only \texttt{7.79744} & Incorrect: only \texttt{7.797} & Multi answer rows need all entries, not just the last value \\
24 & MCQ improper integral & Incorrect & Correct: \texttt{A} & Correct: \texttt{A} & Finalization gives a clean option letter \\
436 & Free response parabola equation & Correct & Correct: \texttt{y\string^2/32} & Incorrect: \texttt{32*x} & Reranking can overwrite a correct answer \\
786 & MCQ integer equation & Incorrect & Incorrect: \texttt{F} & Correct: \texttt{H} & Reranking can recover from a bad candidate \\
\bottomrule
\end{tabular}
}
\end{table}

\section{Discussion}
\label{sec:discussion}

\paragraph{Achievements.}
The LoRA finalizer is the best current default. It improves the held out score from 14/40 to 25/40 and raises the public weighted estimate from 0.4001 to 0.6501. Method 2 ties it in raw count, but Method 1 is simpler and does better on free response rows, which make up most of the public set.

\paragraph{Bottlenecks and Mitigation.}
The largest bottleneck is answer reliability. The model can reason for many tokens and still lose the row because the final answer is incomplete or malformed. The finalizer helps by forcing a concise boxed answer. The remaining weak spot is multi answer free response rows. The id 20 example shows that the finalizer can keep only the last statistic and drop the table entries. Reranking helps some conflicts, but it is not reliable enough to use everywhere.

\paragraph{Plans Forward.}
The next submission should use Method 1 as the default pipeline. Method 2 is worth using selectively on rows where the 1024 and 2048 token answers disagree, especially for multiple choice rows. The most important improvement is targeted handling for multi answer free response rows. GRPO should only be revisited if we can add denser rewards or more training data, since the current outcome reward did not beat supervised finalization.

\paragraph{Team Responsibilities.}
Luis Xu and Wenhao Ren jointly built the Qwen inference and evaluation workflow. Wenhao focused on data inspection, stratified evaluation design, result aggregation, and report organization. Luis focused on pipeline execution, model runs, and comparison experiments. Both team members reviewed the results and chose the final method framing.

\paragraph{Timeline.}
The remaining work is focused on submission readiness and final report polish:
\begin{itemize}
  \item Week 6: Use Method 1 as the main pipeline and keep Method 2 as a selective reranking option.
  \item Week 7: Add targeted handling for multi answer rows and rerun the 40 row held out set.
  \item Week 8: Prepare a Kaggle submission and compare it with the held out results.
  \item Week 9: Add final qualitative examples and a clearer error taxonomy.
  \item Week 10: Finish the report, code release, and presentation materials.
\end{itemize}

\section*{LLM Usage}

We used LLM tools for coding help, experiment organization, debugging, result summaries, and drafting. All experiments were run locally with the provided data, the allowed Qwen model, and the competition judge. We reviewed the text and take responsibility for the claims and results.

\begin{thebibliography}{9}

\bibitem{wei2022cot}
Jason Wei, Xuezhi Wang, Dale Schuurmans, Maarten Bosma, Brian Ichter, Fei Xia, Ed Chi, Quoc Le, and Denny Zhou.
\newblock Chain of thought prompting elicits reasoning in large language models.
\newblock \emph{NeurIPS}, 2022.
\newblock \url{https://arxiv.org/abs/2201.11903}

\bibitem{wang2022selfconsistency}
Xuezhi Wang, Jason Wei, Dale Schuurmans, Quoc Le, Ed Chi, Sharan Narang, Aakanksha Chowdhery, and Denny Zhou.
\newblock Self consistency improves chain of thought reasoning in language models.
\newblock arXiv preprint arXiv:2203.11171, 2022.
\newblock \url{https://arxiv.org/abs/2203.11171}

\bibitem{hu2021lora}
Edward J. Hu, Yelong Shen, Phillip Wallis, Zeyuan Allen Zhu, Yuanzhi Li, Shean Wang, Lu Wang, and Weizhu Chen.
\newblock LoRA: Low rank adaptation of large language models.
\newblock arXiv preprint arXiv:2106.09685, 2021.
\newblock \url{https://arxiv.org/abs/2106.09685}

\bibitem{shao2024deepseekmath}
Zhihong Shao, Peiyi Wang, Qihao Zhu, Runxin Xu, Junxiao Song, Xiao Bi, Haowei Zhang, Mingchuan Zhang, Y. K. Li, Y. Wu, and Daya Guo.
\newblock DeepSeekMath: Pushing the limits of mathematical reasoning in open language models.
\newblock arXiv preprint arXiv:2402.03300, 2024.
\newblock \url{https://arxiv.org/abs/2402.03300}

\end{thebibliography}

\end{document}
